  \documentclass[xcolor=table]{beamer}
\usepackage{beamerthemesplit}
\usepackage{wrapfig}
\usetheme{SPbGU}
\usepackage{pdfpages}
\usepackage{amsmath}
\usepackage{cmap}
\usepackage[T2A]{fontenc}
\usepackage[utf8]{inputenc}
\usepackage[english]{babel}
\usepackage{indentfirst}
\usepackage{amsmath}
\usepackage{tikz}
\usepackage{multirow}
\usepackage[noend]{algpseudocode}
\usepackage{algorithm}
\usepackage{algorithmicx}
\usepackage{fancyvrb}
\usetikzlibrary{shapes,arrows}
%usepackage{fancyvrb}
%\usepackage{minted}
%\usepackage{verbments}


\beamertemplatenavigationsymbolsempty

\title[Результаты группы за 2019 год]{Результаты группы за 2019 год}
\institute[СПбГУ]{
JetBrains Research, Programming Languages and Tools Lab  \\
Санкт-Петербургский Государственный Университет
}

\author[Семён Григорьев]{Семён Григорьев}

\date{14.12.2019}

\begin{document}
{
\begin{frame}[fragile]
  \begin{tabular}{p{2.0cm} p{7.5cm} p{1cm}}
   \begin{center}
      \includegraphics[height=1.5cm]{pictures/jetbrainsResearch.pdf}
    \end{center}
    &
    \begin{center}
      \includegraphics[height=1.5cm]{pictures/YC_logo.pdf}
    \end{center}
    &
    \begin{center}
      \includegraphics[height=1.5cm]{pictures/SPbGU_Logo.png}
    \end{center}
  \end{tabular}
  \titlepage
\end{frame}
}


\begin{frame}[fragile]
  \transwipe[direction=90]
  \frametitle{Области интересов}
\begin{itemize}
      \item Теория формальных языков
      \item Алгоритмы синтаксического анализа
      \item Применение теории формальных языков и синтаксического анализа для решения прикладных задач
      \begin{itemize}
        \item Статический анализ кода
        \item Анализ графовых баз даных
        \item Биоинформатика
        \item \textit{Поиск новых областей}
      \end{itemize}

\end{itemize}

\end{frame}

\begin{frame}[fragile]
  \transwipe[direction=90]
  \frametitle{Состав}
\begin{itemize}
      \item Руководитель: Семён Григорьев
      \item Аспиранты
      \begin{itemize}
        \item Рустам Азимов
        \item Екатерина Шеметова
      \end{itemize}
      \item Магистры
      \begin{itemize}
        \item Юлия Сусанина
        \item Полина Лунина
      \end{itemize}
      \item Студенты
      \begin{itemize}
        \item Максим Шамрай
      \end{itemize}
      \item Формально пока не входят в группу, но работают с нами (курсовые, дипломы, семестровые проекты)
      \begin{itemize}
        \item Арсений Терехов
        \item Дмитрий Кутленков
        \item Артём Хорошев
        \item Илья Эпельбаум
        \item Влада Погожелькая
        \item Вадим Абзалов
        \item Анна Кудрина, Михаил Николюкин, Олег Кулдышев, Максим Юрченков, Анжелина Суханова, Тимур Зиннатулин, \dots
      \end{itemize}

\end{itemize}

\end{frame}


\begin{frame}[fragile]
  \transwipe[direction=90]
  \frametitle{Научные конференции}
\begin{itemize}

      \item \textbf{SIGMOD-2019} (доклад + постер)
      \begin{itemize}
        \item Никита Мишин, Ярослав Соколов, Егор Спирин, Владимир Кутуев, Егор Немчинов, Сергей Горбатюк, \emph{Семён Григорьев}. Evaluation of the Context-Free Path Querying Algorithm Based on Matrix Multiplication (GRADES-NDA)
        \item Публикация: Scopus
      \end{itemize}

      \item \textbf{WoLLIC-2019}
      \begin{itemize}
        \item Сергей Божко, Лейла Хатбуллина, \emph{Семён Григорьев}. Bar-Hillel Theorem Mechanization in Coq
        \item Публикация: LNCS, Scopus, WoS
      \end{itemize}

      \item \textbf{SECR-2019}
      \begin{itemize}
         \item \emph{Илья Ножкин}, Семнё Григорьев. CFL-Reachability Based Framework for Interprocedural Static Code Analysis Development
         \item Bertrand Meyer award
         \item Публикация: Scopus
      \end{itemize}

    \end{itemize}
    \end{frame}


\begin{frame}[fragile]
  \transwipe[direction=90]
  \frametitle{Научные конференции}
      \begin{itemize}

      \item \textbf{BIOINFORMATICS-2019}
      \begin{itemize}
        \item \emph{Семён Григорьев}, Полина Лунина. The Composition of Dense Neural Networks and Formal Grammars for Secondary Structure Analysis
        \item Публикация: Scopus
      \end{itemize}

      \item \textbf{CIBB-2019}
      \begin{itemize}
        \item \emph{Полина Лунина}, Семён Григорьев. On Secondary Structure Analysis by Using Formal Grammars and Artificial Neural Networks
        \item \emph{Юлия Сусанина}, Семён Григорьев. Modification of Valiant’s Parsing Algorithm for String-Searching Problem
        \item Публикация: ждём
      \end{itemize}

      \item \textbf{BIATA-2019} (постер)
      \begin{itemize}
         \item \emph{Семён Григорьев}, Полина Лунина. Improved Architecture of Artificial Neural Network for Secondary Structure Analysis
         \item Публикация: BMC Bioinformatics, Scopus
      \end{itemize}

\end{itemize}
\end{frame}

\begin{frame}[fragile]
  \transwipe[direction=90]
  \frametitle{Доклады}
\begin{itemize}

      \item Открытая лекция в CSC ``Теория формальных языков на практике''

      \item Доклад на семинаре группы LINKS (Inria, Лилль, Франция)

\end{itemize}
\end{frame}


\begin{frame}[fragile]
  \transwipe[direction=90]
  \frametitle{Публикации}
\begin{itemize}
      \item Шеметова Е.Н., Григорьев С.В. Задача поиска путей в ациклических графах с ограничениями в терминах булевых грамматик. Труды Института системного программирования РАН
\end{itemize}
\end{frame}

\begin{frame}[fragile]
  \transwipe[direction=90]
  \frametitle{Сотрудничество}
\begin{itemize}
      \item Грант РНФ под руководством Александра Охотина
      \begin{itemize}
        \item Семён Григорьев, Рустам Азимов, Екатерина Шеметова
      \end{itemize}
      \item INRIA LINKS (\url{https://team.inria.fr/links/})
      \begin{itemize}
         \item Стажировка Рустама Азимова (апрель 2019)
         \item Поездка Семёна Григорьева (апрель 2019)
         \item Планируется поездка Рустама Азимова
      \end{itemize}
      \pause
      \item Ревью статьи в журнал Science of Computer Programming
      \item Член программного коммитета секции по алгебраическим методам для анализа РНК при конференции CIBB

\end{itemize}
\end{frame}

\begin{frame}[fragile]
  \transwipe[direction=90]
  \frametitle{Образовательная деятельность}
\begin{itemize}
      \item Курс (лекции/семинары) по теории формальных языков
      \item Курс по теории графов на математико-механическом факультете СПбГУ
      \item Семинар по теории формальных языков
      \item Руководство курсовыми, ВКР, магистерскими, семестровыми проектами и т.д. в CSC, Институте биоинформатики, СПбГУ и других ВУЗах
\end{itemize}
\end{frame}

\begin{frame}[fragile]
  \transwipe[direction=90]
  \frametitle{Гранты}
\begin{itemize}
      \item РФФИ ``Конкурс на лучшие проекты фундаментальных научных исследований, выполняемые молодыми учеными, обучающимися в аспирантуре''
      \begin{itemize} {\small
        \item ``Поиск путей с ограничениями в терминах формальных языков''
        \item Семён Григорьев (руководитель), Рустам Азимов (исполнитель, аспирант)
        \item Поддержана: 3 года, 1200 т.р.}
      \end{itemize}

      \item РНФ ``Проведение инициативных исследований молодыми учеными''
      \begin{itemize}{\small
        \item ``Использование формальных грамматик и искусственных нейронных сетей для анализа вторичной структуры геномных и протеомных последовательностей''
        \item Семён Григорьев
        \item Не поддержана}
      \end{itemize}

      \item NVIDIA Gaduate Fellowship
      \begin{itemize}{\small
        \item ``An Efficient GPU-based Path Querying of the Graph Databases''
        \item Рустам Азимов
        \item Не поддержана}
      \end{itemize}

      \item РНФ ``Проведение фундаментальных научных исследований и поисковых научных исследований отдельными научными группами''
      \begin{itemize}{\small
        \item ``Теория формальных языков и алгоритмы синтаксического анализа для анализа граф-структурированных данных''
        \item Семён Григорьев, Екатерина Вербицкая, Даниил Березун, Рустам Азимов, Екатерина Шеметова, Юлия Сусанина, Арсений Терехов, Илья Балашов, Никита Мишин
        \item Ждём до 1 апреля}
      \end{itemize}
\end{itemize}
\end{frame}

\begin{frame}[fragile]
  \transwipe[direction=90]
  \frametitle{Планируемые публикации}
\begin{itemize}
      \item !!!
      \item !!!
      \item !!!
\end{itemize}
\end{frame}

\begin{frame}[fragile]
  \transwipe[direction=90]
  \frametitle{Ведущиеся работы}
\begin{itemize}
      \item Исследование матричных алгоритмов решения задачи контекстно-свободной достижимости
      \item Субкубический алгоритм для контекстно-свободной достижимости
      \item Применение синтаксического анализа для обработки биологических последовательностей
      \item Построение групп по формальному языку, использование теоретико-групповых подходов к задачам теории формальных языков
      \item Применеие специализации и других техник частичных вычислений для оптимизации процедур выполнения запросов к большим данным
      \begin{itemize}
        \item Совместно с группой Даниила Березуна
      \end{itemize}
\end{itemize}
\end{frame}

\end{document}
